\subsubsection{The self-duality property. Tube type Hermitian symmetric spaces}
By definition, a Hermitian symmetric space $M=G_{\mathbb R}/K_{\mathbb R}$ is of \emph{tube type} if it is biholomorphically equivalent to a domain $F\oplus iC$, where $C\subset F$ is a proper open cone in the real vector space $F$. The simplest example is the upper half-plane, where $F=\mathbb R$ and $C=\mathbb R^*_+$. For irreducible symmetric spaces, this happens when $G_{\mathbb R}$ is isogenous to $\mathrm{SU}(p,p)$, $\mathrm{Spin}(2,m)$, $\mathrm{Sp}(2m,\mathbb R)$, $\mathrm{SO}^{\star}(2m)$ for $m$ even, and $\mathrm E_{7(-25)}$.

We introduce the following definition.
\begin{definition}
A sequence $(\gamma_1,\ldots,\gamma_r)$ of roots will be called \emph{admissible} if all the roots $\gamma_i$ are long and bigger than $\zeta$, and if they are pairwise strongly orthogonal.
\end{definition}
Note that the dominant orthogonal sequences $(\alpha_1,\ldots,\alpha_r)$ are admissible. We can now state:
\begin{proposition}
The following assertions are equivalent:
\begin{enumerate}[label=(\arabic*)]
\item $M$ has tube type
\item For all admissible sequences $\gamma_1,\ldots,\gamma_{r_M}$, we have $\gamma_1+\cdots+\gamma_{r_M}=2\varpi$
\item $\alpha_1+\cdots+\alpha_{r_M}=2\varpi$
\item There exists an admissible sequence $\gamma_1,\ldots,\gamma_{r_M}$ such that $\gamma_1+\cdots+\gamma_{r_M}=2\varpi$
\item $\alpha_1^\jepPubCoroot+\cdots+\alpha_{r_M}^\jepPubCoroot=z$
\item $\jepPubBbb{E}$ is self-dual: $\jepPubBbb{E}\simeq\jepPubBbb{E}^{\star}$ as $G$-modules
\item For all $i$, $\jepPubBbb{E}_i\simeq\jepPubBbb{E}_{r_M-i}^{\star}$ as $K$-modules
\item $\dim\jepPubBbb{E}_{r_M}=1$
\item $\alpha_{r_M}=\zeta$
\end{enumerate}
\end{proposition}
\begin{proof}
Glancing at Table~1, one can readily check that $G_{\mathbb R}/K_{\mathbb R}$ is of tube type if and only if $z=\alpha_1^\jepPubCoroot+\cdots+\alpha_{r_M}^\jepPubCoroot$, and that this occurs also exactly when $\alpha_{r_M}=\zeta$. Hence~(1),~(5) and~(9) are equivalent.

The equality $z=a_1^\jepPubCoroot+\cdots+\alpha_{r_M}^\jepPubCoroot$ is equivalent to $\langle\alpha_1^\jepPubCoroot+\cdots+\alpha_{r_M}^\jepPubCoroot,\beta\rangle=2\delta_{\beta,\zeta}$ for all $\beta\in\Pi$ because $\zeta$ is the only noncompact simple root and it has the same length as each $\alpha_i$ (they are all long). Since all the roots $\alpha_i$ have the same length, for any root $\beta$, the equation $\langle\alpha_1^\jepPubCoroot+\cdots+\alpha_{r_M}^\jepPubCoroot,\beta\rangle=0$ is equivalent to $\langle\alpha_1+\cdots+\alpha_{r_M},\beta^\jepPubCoroot\rangle=0$. It follows that~(5) and~(3) are equivalent.

We now prove that items (2), (3), and (4) are equivalent. In fact, we have the trivial implications $\text{(2)}\Rightarrow\text{(3)}\Rightarrow\text{(4)}$. Moreover, if~(4) holds, then let $\gamma_1,\ldots,\gamma_{r_M}$ be an admissible sequence such that $\gamma_1+\cdots+\gamma_{r_M}=2\varpi$. By \cite[Th.~2.1]{jep93RRS92}, any admissible sequence can be obtained applying to this sequence an element $w$ of the subgroup of the Weyl group that fixes $\varpi$. Condition~(2) follows from the equality $w(\gamma_1)+\cdots+w(\gamma_{r_M})=w(2\varpi)=2\varpi$.

Since the weights of $\jepPubBbb{E}^{\star}$ are the opposite of the weights of $\jepPubBbb{E}$, the highest weight of $\jepPubBbb{E}^{\star}$ is $-\mu_0$. Therefore,~(6) holds if and only if $\mu_0=-\varpi$, which is equivalent to~(3) by Proposition~2.19~(4).

Assuming~(6) and using the decomposition into $K$-modules $\jepPubBbb{E}=\jepPubOplus_{i=0}^{r_M}\jepPubBbb{E}_i$, we deduce that $\jepPubBbb{E}_{r_M-i}\simeq\jepPubBbb{E}_i^{\star}$ as $K$-modules. Conversely, given~(7), we obviously have~(8). Assuming now~(8), we deduce that $\jepPubBbb{E}_{r_M}=\jepPubBbb{E}_{\mu_0}$. Since $\jepPubBbb{E}_{r_M}$ is a $K$-module of dimension $1$ in this case, it follows that $\mu_0$ vanishes on all the roots of $\mathfrak k$, so $\mu_0$ is an integral multiple of $\varpi$. It follows that $\mu_0=-\varpi$ and we get~(3) as above.
\end{proof}
As we just said, the $G$-module $\jepPubBbb{E}$ is not always self-dual. However, the subspace $\jepPubBbb{V}^r$ is always a self-dual module under a subgroup $\overline L_r$ of $G$ that we now define.

Recall that the Lie algebra $\mathfrak l_r$ is by definition generated by the root spaces of the roots $\alpha$ such that $\langle h_r,\alpha\rangle=\langle z,\alpha\rangle$, and that $\mathfrak l_r^\pm$ is the subalgebra generated by the root spaces of the roots $\alpha$ such that $\langle h_r,\alpha\rangle=\langle z,\alpha\rangle=\pm2$. We define $\overline{\mathfrak l}_r$ as the Lie subalgebra of $\mathfrak g$ generated by the spaces $\mathfrak l_r^+$ and $\mathfrak l_r^-$. Two particular cases are easy:
\begin{fact}
If $r=1$, then $\mathfrak{sl}_2\simeq\overline{\mathfrak l}_r\subset\mathfrak g$ is the Lie subalgebra corresponding to the highest root $\Theta$. If $M$ has tube type and $r=r_M$, then $\overline{\mathfrak l}_r=\mathfrak g$.
\end{fact}
\begin{proof}
In case $r=1$, then $h_r=\Theta^\jepPubCoroot$, so $\langle\alpha,h_r\rangle=2$ implies $\alpha=\Theta$. Thus $\mathfrak l_r^+=\mathfrak g_\Theta$ in this case. When $M$ has tube type and $r=r_M$, we have $h_r=z$ (Proposition~2.35~(5)), so $\mathfrak l_r^+=\mathfrak m^+$ in this case.
\end{proof}
We now describe the Lie algebra $\overline{\mathfrak l}_r$ in general. To this end, we introduce a definition.
\begin{definition}
Let $R$ be a (possibly reducible) root system, let $\Pi$ be a basis of $R$, and $\varpi$ be a dominant weight. Given a simple root $\beta$, we say that $\beta$ is \emph{connected to $\varpi$ in $R$} if there is a simple root $\gamma$ in the same connected component as $\beta$ in the Dynkin diagram of $R$ and such that $\langle\varpi,\gamma^\jepPubCoroot\rangle>0$.
\end{definition}
We denote by $\Pi_{h_r=0}$ the set of simple roots $\beta\in\Pi$ such that $\langle\beta,h_r\rangle=0$.
\begin{lemma}
The Lie algebra $\overline{\mathfrak l}_r$ is simple. The lowest root $-\Theta$ together with the simple roots which are connected to $\Theta$ in $\Pi_{h_r=0}$ form a basis of the root system of $\overline{\mathfrak l}_r$.
\end{lemma}
\begin{proof}
Recall that $\zeta$ denotes the only noncompact simple root (see Notation~2.5). Looking at Table~1, we see that the root $\zeta$ cannot be connected to $\Theta$ in $\Pi_{h_r=0}$. Thus, if $\gamma$ is connected to $\Theta$ in $\Pi_{h_r=0}$, it satisfies $\langle\gamma,z\rangle=0$ and, of course, $\langle\gamma,h_r\rangle=0$. Let us assume that $\langle\Theta,\gamma^\jepPubCoroot\rangle>0$. Then $\Theta-\gamma$ is a root and it satisfies $\langle h_r,\Theta-\gamma\rangle=\langle z,\Theta-\gamma\rangle=2$, so $\mathfrak g_{\Theta-\gamma}\subset\overline{\mathfrak l}_r$, and so $\mathfrak g_\gamma\subset\overline{\mathfrak l}_r$. By induction, the root spaces of all simple roots connected to $\Theta$ and their opposite are in $\overline{\mathfrak l}_r$, and, together with $\gamma_{-\Theta}$ and $\gamma_\Theta$, they generate $\overline{\mathfrak l}_r$.

In general, the Dynkin diagram of a Lie algebra has vertices corresponding to the simple roots of the Lie algebra. In our case, given the above description of a basis of $\overline{\mathfrak l}_r$, the set of vertices of its Dynkin diagram is the union of some connected components of $\Pi_{h_r=0}$ and the set $\{-\Theta\}$. Moreover, all the connected components of $\Pi_{h_r=0}$ which occur are connected to $-\Theta$. Thus the Dynkin diagram of $\overline{\mathfrak l}_r$ is connected and $\overline{\mathfrak l}_r$ is simple.
\end{proof}
In Table~3, we indicate in each case which root system for $\overline{\mathfrak l}_r$ is obtained (the cases where $r=1$ or $M$ has tube type and $r=r_M$ are not represented, see Fact~2.36).
\clearpage

\begin{table}[htbp]\centering\normalsize
\setlength{\tabcolsep}{5pt}\renewcommand{\arraystretch}{1.4}
\begin{tabular}{|c|c|c|}\hline
$(G,\varpi)$ & $\overline{\mathfrak l}_r$ & $R(\overline{\mathfrak l}_r)$\\\hline
$(A_{n-1},\varpi_\ell)$ & \jepGraph{0/0/0/{\beta_{r-1}}/{}/0,1/1.55/0/{\beta_2}/{}/0,2/2.55/0/{\beta_1}/{}/0,3/3.55/0/{-\Theta}/{}/1,4/4.55/0/{\beta_n}/{}/0,5/5.55/0/{\beta_{n-1}}/{}/0,6/7.1/0/{\beta_{n-r+1}}/{}/0}{1/2,2/3,3/4,4/5}{0/1,5/6}{}{} & $A_{2r-1}$\rule[-15pt]{0pt}{46pt}\\\hline
$(C_n,\varpi_n)$ & \jepGraph{0/0/0/{\beta_{r-1}}/{}/0,1/1.55/0/{\beta_1}/{}/0,2/2.55/0/{-\Theta}/{}/1}{}{0/1}{1/2/-1}{} & $C_r$\rule[-15pt]{0pt}{46pt}\\\hline
$(D_n,\varpi_n)$ & \jepGraph{0/0/0/{\beta_{2(r-1)}}/{}/0,1/1.55/0/{\beta_3}/{}/0,2/2.55/0/{\beta_2}/{}/0,3/3.05/0.8/{\beta_1}/{}/0,4/3.05/-0.8/{-\Theta}/{}/1}{1/2,2/3,2/4}{0/1}{}{} & $D_{2r-1}$\rule[-15pt]{0pt}{46pt}\\\hline
$(E_6,\varpi_1)$, $r=2$ & \jepGraph{0/0/0/{-\Theta}/{}/1,1/1/0/{\beta_2}/{}/0,2/2/0/{\beta_3}/{}/0,3/2.5/0.8/{\beta_1}/{}/0,4/2.5/-0.8/{\beta_4}/{}/0}{0/1,1/2,2/3,2/4}{}{}{} & $D_5$\rule[-15pt]{0pt}{46pt}\\\hline
$(E_7,\varpi_7)$, $r=2$ & \jepGraph{0/0/0/{-\Theta}/{}/1,1/1/0/{\beta_1}/{}/0,2/2/0/{\beta_3}/{}/0,3/3/0/{\beta_4}/{}/0,4/3.5/0.8/{\beta_5}/{}/0,5/3.5/-0.8/{\beta_2}/{}/0}{0/1,1/2,2/3,3/4,3/5}{}{}{} & $D_6$\rule[-15pt]{0pt}{46pt}\\\hline
\end{tabular}
\caption{\small The Lie algebra $\overline{\mathfrak l}_r$ and its root system (The cases where $r=1$ or $M$ has tube type and $r=r_M$ are not represented.)}
\label{jep93table3}
\end{table}
\clearpage

We consider the root system $R_{\overline{\mathfrak l}_r}$ with its basis $(-\Theta,\beta_1,\ldots,\beta_k)$ given by Lemma~2.38. We have
\begin{lemma}
The simple root $-\Theta$ in $R(\overline{\mathfrak l}_r)$ is cominuscule.
\end{lemma}
In other words, the lemma says that any root in $R(\overline{\mathfrak l}_r)$ has coefficient at most one on the simple root $-\Theta$. Observe that, as a consequence, $-\Theta$ defines a Hermitian noncompact real form $\overline L_{r,\mathbb R}$ of $\overline L_r$. The element $y=y^r$ can be considered as an element of $\mathfrak l_r^-$.
\begin{proof}
Since any root in $R(\overline{\mathfrak l}_r)$ can be obtained from $\Theta$ applying a suitable element of the Weyl group of $\overline{\mathfrak l}_r$, it is enough to show the implication
\[
|n_{-\Theta}(\alpha)|\leqslant1\Longrightarrow|n_{-\Theta}(s_\gamma(\alpha))|\leqslant1,
\]
for any simple root $\gamma$ in $R(\overline{\mathfrak l}_r)$. To prove this, we assume without loss of generality that $n_{-\Theta}(\alpha)\geqslant0$. If $\gamma$ is one of the simple roots $\beta_j$, then $n_{-\Theta}(\alpha)=n_{-\Theta}(s_\gamma(\alpha))$, so the implication clearly holds. If $\gamma=-\Theta$, we consider two cases. If $n_{-\Theta}(\alpha)=1$, then $\alpha$ is a negative root. We have $0\leqslant\langle\alpha,-\Theta^\jepPubCoroot\rangle\leqslant2$, thus $-1\leqslant n_{-\Theta}(s_\gamma(\alpha))\leqslant1$. If $n_{-\Theta}(\alpha)=0$, then $\alpha\neq-\Theta$. Since $\Theta$ is long, it follows that $|\langle\alpha,-\Theta^\jepPubCoroot\rangle|\leqslant1$. Thus $|n_{-\Theta}(s_\gamma(\alpha))|\leqslant1$, as expected.
\end{proof}
As announced, we have the following, which strengthens Lemma~2.27.
\begin{proposition}
The $\overline L_r$-module generated by $\jepPubBbb{E}_\varpi$ is $\jepPubBbb{V}^r$. It is self-dual.
\end{proposition}
\begin{proof}
First, we have $\overline{\mathfrak l}_r\subset\mathfrak l_r$. This follows from the definition of $\overline{\mathfrak l}_r$ and the inclusions $\mathfrak l_r^+\subset\mathfrak l_r,\ \mathfrak l_r^-\subset\mathfrak l_r$.

From the inclusion $\overline{\mathfrak l}_r\subset\mathfrak l_r$ and Proposition~2.33(1), it follows that the $\overline L_r$-submodule generated by $\jepPubBbb{E}_\varpi$ is included in $\jepPubBbb{V}^r$. By Lemma~2.31, $\jepPubBbb{V}^r$ is included in the $\mathfrak l_r^-$-submodule generated by $\jepPubBbb{E}_\varpi$. Since $\mathfrak l_r^-\subset\overline{\mathfrak l}_r$ by definition, we deduce that the $\overline L_r$-module generated by $\jepPubBbb{E}_\varpi$ is exactly $\jepPubBbb{V}^r$.

Now we prove that $\jepPubBbb{V}^r$, as an $\overline L_r$-module, is self-dual. The highest weight of $\jepPubBbb{V}^r$ as an $\overline L_r$-module is the restriction of $\varpi$, which is the fundamental weight $\varpi_{-\Theta}$ corresponding to the simple root $-\Theta$ in the root system of $\overline{\mathfrak l}_r$. The sequence $\alpha_1,\ldots,\alpha_r$ is a sequence of long strongly orthogonal roots. They are all bigger than $-\Theta$. In fact, if $\alpha_i$ was not bigger than $-\Theta$, then it would be a linear combination of the roots $\beta_j$, so it would satisfy $n_\zeta(\alpha_i)=0$, contradicting the definition of the sequence $\alpha_1,\ldots,\alpha_r$.

In view of~(4) in Proposition~2.35, it is enough to prove that $\sum\alpha_i$ and $2\varpi_{-\Theta}$ have the same pairing with the simple roots of $R_{\overline{\mathfrak l}_r}$. For $-\Theta$, we have $\langle\sum\alpha_i^\jepPubCoroot,-\Theta^\jepPubCoroot\rangle=-2$ and $\langle2\varpi_{-\Theta},-\Theta^\jepPubCoroot\rangle=-2$. For $\beta$ a simple root connected to $\Theta$ in $\Pi_{h_r=0}$, we have $\langle\sum\alpha_i,\beta^\jepPubCoroot\rangle=0$ because $\langle\beta,\sum\alpha_i^\jepPubCoroot\rangle=0$. We have $\langle\varpi_{-\Theta},\beta^\jepPubCoroot\rangle=0$. This proves that $\jepPubBbb{V}^r$ is self-dual.
\end{proof}
From the equivalences in Proposition~2.35 and Proposition~2.40, and recalling that $\overline L_{r,\mathbb R}$ is a Hermitian noncompact real form of $\overline L_r$, it follows that the symmetric space $M^r:=\overline L_{r,\mathbb R}/(\overline L_{r,\mathbb R}\cap K_{\mathbb R})$ is a tube type Hermitian symmetric space of rank $r$.
\begin{remark}
This construction is related, but not identical, to the theory of boundary components of bounded symmetric domains (see e.g. \cite{jep93Wol72,jep93PS69,jep93Sat80,jep93AMRT10}). The topological boundary of the symmetric space $M$ inside its compact dual $\check M$ is a union of $r_M$ $G_{\mathbb R}$-orbits, and each $G_{\mathbb R}$-orbit is a union of so-called boundary components, cf. \cite[Def.~3.2 p.~127]{jep93AMRT10}. The components are isomorphic to Hermitian symmetric subspaces $M_\Lambda$ of $M$, where $\Lambda$ is a subset of the set of noncompact strongly orthogonal roots, see \cite[Th.~3.3 p.~127]{jep93AMRT10} for example. If we take our maximal dominant orthogonal sequence $(\alpha_1,\ldots,\alpha_{r_M})$ as the set of strongly orthogonal roots, and $\Lambda=\{\alpha_1,\ldots,\alpha_r\}$ as a subset, boundary components corresponding to this subset are symmetric spaces of rank $r$, as is the subspace $M^r$, but they are not of tube type unless $M$ already is. The Lie algebra of the stabilizer of $M_\Lambda$ in $G_{\mathbb R}$ is $\mathfrak q_\Lambda\cap\mathfrak g_{\mathbb R}$, where $\mathfrak q_\Lambda$ is the parabolic subalgebra of $\mathfrak g$ corresponding to the coweight $h_{r_M}-h_r$, whereas we recall that the parabolic subalgebra $\mathfrak q_r$ we are considering here corresponds to the coweight $z-h_r$. This shows that $\mathfrak q_r=\mathfrak q_\Lambda$ if and only if $M$ is of tube type.
\end{remark}
\clearpage

\begin{example}
We give an example of this construction. We assume that we are in the first row of Table~1, namely that $G$ has type $A_{p+q-1}$ for some positive integers $p\leqslant q$ and that $\varpi=\varpi_p$. The group $G_{\mathbb R}$ is therefore $\operatorname{SU}(p,q)$ acting on $\mathbb C^{p+q}$ preserving a Hermitian form $h$ of signature $(p,q)$ given by the diagonal matrix $\operatorname{diag}(-1,\ldots,-1,1,\ldots,1)$ in the canonical basis $(e_1,\ldots,e_p,e_{p+1},\ldots,e_{p+q})$.

Let $y\in\mathfrak m^-$ be an element of rank $r$. We have a natural decomposition $\mathbb C^{p+q}=A\oplus N\oplus B\oplus I$, with $A$, resp. $N,B,I$ generated by $(e_1,\ldots,e_r)$, resp. $(e_{r+1},\ldots,e_p)$, $(e_{p+1},\ldots,e_{p+q-r})$, $(e_{p+q-r+1},\ldots,e_{p+q})$, so that $y$ is given by the block matrix
\[
\begin{pmatrix}0&0&0&0\\0&0&0&0\\0&0&0&0\\ *&0&0&0\end{pmatrix}.
\]
The element $h_r\in\mathfrak t$ acts on the Lie algebra $\mathfrak{sl}(p+q,\mathbb C)$ of $G$, and we represent the weights of its action as the matrix
\[
\begin{pmatrix}0&1&1&2\\-1&0&0&1\\-1&0&0&1\\-2&-1&-1&0\end{pmatrix},
\]
where the blocks correspond to the decomposition $A\oplus N\oplus B\oplus I$ of $\mathbb C^{p+q}$. Similarly, the weights of the action of the central element $z$ are given as follows
\[
\begin{pmatrix}0&0&2&2\\0&0&2&2\\-2&-2&0&0\\-2&-2&0&0\end{pmatrix}.
\]
Thus the Lie algebra $\mathfrak q_r$, $\mathfrak l_r$ and $\mathfrak h_r$ are respectively
\[
\begin{pmatrix}*&0&*&*\\{*}&*&*&*\\0&0&*&0\\{*}&0&*&*\end{pmatrix},\quad
\begin{pmatrix}*&0&0&*\\0&*&0&0\\0&0&*&0\\{*}&0&0&*\end{pmatrix}
\quad\text{and}\quad
\begin{pmatrix}*&0&0&0\\0&*&0&0\\0&0&*&0\\0&0&0&*\end{pmatrix}.
\]
We see that $Q_r$ is exactly the stabilizer of the flag $(N\subset N\oplus A\oplus I)$, that the subgroup $L_r$ of $G$ corresponding to $\mathfrak l_r$ is $\operatorname S(\operatorname{GL}(N)\times\operatorname{GL}(A\oplus I)\times\operatorname{GL}(B))$ and that $H_r$ is the block diagonal group $\operatorname S(\operatorname{GL}(N)\times\operatorname{GL}(A)\times\operatorname{GL}(I)\times\operatorname{GL}(B))$. Finally the intersections of $G_{\mathbb R}=\operatorname{SU}(p,q)$ with these two latter groups are isomorphic to $\operatorname S(\operatorname U(p-r)\times\operatorname U(r,r)\times\operatorname U(q-r))$ and $\operatorname S(\operatorname U(p-r)\times\operatorname U(r)\times\operatorname U(r)\times\operatorname U(q-r))$ respectively, so that $\overline L_{r,\mathbb R}\simeq\operatorname{SU}(r,r)$.

On the other hand, we have $\jepPubBbb{E}=\wedge^p(N\oplus A\oplus I\oplus B)$ and it is easy to check that $\jepPubBbb{V}^r=\wedge^{p-r}N\otimes\wedge^r(A\oplus I)$, which illustrates that the stabilizer of $\jepPubBbb{V}^r$ is $Q_r$.
\end{example}
